% ============================================================================
%  CUMCM 2026 本科组 A/B/C 题 论文模板（电子版）
%  依据《全国大学生数学建模竞赛论文格式规范（2026年修订稿）》编写
%  章节骨架与篇幅参数来自 44 篇官方展示论文（一等奖级）的实测分位
%
%  编译：xelatex main.tex （跑两遍，第二遍生成交叉引用）
%        本模板已实测编译通过：TeX Live 2026 + xelatex，5 页 / 0 错误 / 0 警告
%
%  字体：必须用 fontset=windows（SimSun/SimHei）。
%        ★ 实测教训：fontset=fandol 编译不报错、显示也正常，但 FandolSong
%        以 Identity-H 嵌入且【不带 ToUnicode CMap】，导致 PDF 里的中文
%        无法复制、无法检索、相似度检测可能读不到内容。换 windows 后
%        SimSun/SimHei 的 ToUnicode 正常，中文可提取。
%        无中易字体的机器（Linux/macOS）改 fontset=ubuntu / macnew，
%        换完务必重新验证一次中文能不能从 PDF 里复制出来。
%
%  【格式规范硬约束——改动前务必核对】
%   第三条  摘要专用页，含标题与关键词，不能超过一页；页码从此页起
%   第四条  正文从摘要页之后开始，不要目录，尽量控制在 20 页以内
%   第六条  全文任何地方不得出现姓名、学校、赛区信息
%   第十条  电子版第一页必须是摘要页（不含承诺书、编号页），PDF，≤20MB
%   AI 规定 “AI 工具使用声明”必须放在【参考文献之前】，用官方原句
% ============================================================================
\documentclass[UTF8,12pt,a4paper,fontset=windows]{ctexart}

\usepackage[margin=2.5cm]{geometry}   % 规范第一条：四边至少 2.5cm
\usepackage{amsmath,amssymb,amsthm,bm}
\usepackage{graphicx,float,subcaption}
\usepackage{booktabs,longtable,multirow,array}
\usepackage{listings,xcolor}
\usepackage{caption}
\usepackage{enumitem}
\usepackage[hidelinks,unicode]{hyperref}

% —— 页码：页脚居中，阿拉伯数字，从摘要页起连续（规范第三条）——
\usepackage{fancyhdr}
\pagestyle{fancy}\fancyhf{}\fancyfoot[C]{\thepage}
\renewcommand{\headrulewidth}{0pt}

\setlength{\parindent}{2em}
\linespread{1.3}
\captionsetup{font=small,labelfont=bf}
\numberwithin{equation}{section}

% —— 附录代码环境：行号 + 语法高亮，便于评委阅读 ——
\definecolor{cmt}{RGB}{0,128,0}\definecolor{kwd}{RGB}{0,0,220}\definecolor{str}{RGB}{163,21,21}
\lstset{
  basicstyle=\footnotesize\ttfamily, breaklines=true, breakindent=2em,
  frame=single, numbers=left, numberstyle=\tiny\color{gray},
  commentstyle=\color{cmt}, keywordstyle=\color{kwd}\bfseries,
  stringstyle=\color{str}, showstringspaces=false,
  columns=fullflexible, keepspaces=true, tabsize=4,
  extendedchars=true
}

% —— 符号说明表的便捷命令 ——
\newcommand{\symrow}[3]{$#1$ & #2 & #3 \\}

\title{\bfseries 在此填写论文标题}
\date{}
\author{}

\begin{document}
% ============================================================================
%  第 1 页：摘要专用页（不能超过一页）
%  实测参数（44 篇一等奖级）：中文 848 / 927 / 1010 字（p25/p50/p75），关键词 4 个
%  成分命中率：91% 写出具体数值结果，86% 用“针对问题X”分段，82% 明写模型名
%  【别写模型评价/优缺点】只有 11% 这么做，那是正文倒数第二章的事
%  【别硬凑创新点】只有 5% 提；组委会把“声称创新却无正文证据”列为反模式
% ============================================================================
\maketitle
\thispagestyle{fancy}
\vspace{-3em}

\begin{center}\zihao{4}\bfseries 摘\quad 要\end{center}

% 开场 2–4 句：背景一句（别抄题面）+ 本文建立了什么框架、研究什么问题
本文针对……问题，建立了……模型与……模型相结合的……框架，系统研究了……。

% 每个子问题一段。固定内涵：目标 → 模型 → 算法 → 【具体数值结果】
% 末句给数值时最好带对照基线，例如“较…基线降低 1.8 个百分点”
\textbf{针对问题一}，首先……。通过……，建立……模型，并……。
结果表明：……，模型拟合优度达 $R^2=0.8962$，优于……。

\textbf{针对问题二}，本文以……为目标，构建……优化模型，
利用……算法求解，得到……。分组结果为：……。

\textbf{针对问题三}，……。经交叉验证与灵敏度回验，结论稳健。

\textbf{针对问题四}，……。在……约束下，得到……。

% 稳健性一句话收尾（32% 的一等奖摘要会写，写了是加分）
经灵敏度检验，……对结果的影响较小。

\vspace{1em}
\noindent\textbf{关键词：}\quad 关键词一\quad 关键词二\quad 关键词三\quad 关键词四
% 关键词 4 个左右，务必含【问题域词】而不是全是方法名

\newpage
% ============================================================================
%  正文从第 2 页开始。规范第四条：不要目录，尽量控制在 20 页以内
%  主章节数实测 p50 = 10，子章节 p50 = 31，公式块 p50 = 49，图 p50 = 20，表 p50 = 9
% ============================================================================

\section{问题重述}
\subsection{问题背景}
% 用自己的话概括场景，不要复制题面原文
\subsection{问题提出}
% 按题面任务拆成可操作的子问题，数量与后文章节一一对应

\section{模型假设}
% 每条假设都要有依据，且能在结果之后回检。假设之间不能互相打架
\begin{enumerate}[label=(\arabic*),itemsep=2pt]
  \item 假设……，依据是……；
  \item 假设……，依据是……。
\end{enumerate}

\section{符号说明}
\begin{table}[H]\centering\small
\caption{符号说明}
\begin{tabular}{clc}
\toprule
符号 & 含义 & 单位 \\
\midrule
\symrow{D}{水深}{m}
\symrow{\theta}{换能器开角}{rad}
\bottomrule
\end{tabular}
\end{table}
% 符号不得复用；单位必须写；正文实际用到的符号都要在此出现

\section{问题分析}
% 每个子问题：本质是什么类型的问题 → 难点在哪 → 如何转化为数学问题
%            → 与上游子问题是否复用结果
% 【组委会连续四年批评“不从问题出发，乱套方法”】
% 自检：把这一节换到另一道题上，如果照样成立，就是套路，必须重写

\section{问题一模型的建立与求解}
\subsection{模型建立}
% 机理类：先给解析表达（牛顿定律/微分方程/几何关系），再谈数值
% 数据类：先判断要不要预处理——原始流水数据要清洗；已整理的观测表默认不清洗
\begin{equation}
  y = f(x;\bm{\theta})
\end{equation}

\subsection{求解算法}
% 【必须给出具体算法】2025A 与 2023B 讲评都点名“没有给出具体算法/计算方法不完善”
% 要写：初始化、迭代式、终止条件、步长/容差、求解器与版本
\begin{enumerate}[label=Step \arabic*.,itemsep=2pt]
  \item 初始化：……
  \item 迭代：……
  \item 终止条件：……
\end{enumerate}

\subsection{结果与分析}
% 【五件必交】① 具体数值表 ② 计算方法说明 ③ 合规的结果文件
%             ④ 对结果的分析（不只罗列）⑤ 稳健性/灵敏度
% 2022A 讲评 9 条评阅问题里 5 条是“呈现与分析”，不是建模
\begin{table}[H]\centering\small
\caption{问题一计算结果}
\begin{tabular}{lccc}
\toprule
情形 & 指标一 & 指标二 & 指标三 \\
\midrule
1 & & & \\
\bottomrule
\end{tabular}
\end{table}
% 分析：这些数字的物理/业务含义、是否符合直觉、异常点归因、与基线对比

\begin{figure}[H]\centering
% \includegraphics[width=0.7\textwidth]{fig/q1.pdf}
\caption{问题一结果示意}
\end{figure}

\section{问题二模型的建立与求解}
% 同上结构。命题人的解法节奏几乎都是【从最简情形起步，逐步加约束】，
% 每一步都要说清“上一步为什么不够”
\subsection{模型建立}
\subsection{求解算法}
\subsection{结果与分析}

\section{问题三模型的建立与求解}
\subsection{模型建立}
\subsection{求解算法}
\subsection{结果与分析}

\section{模型的检验与灵敏度分析}
% 参数扰动后结论是否翻转；给出决策/结果的翻转边界
% 极限退化自检：新引入的角度/参数取边界值时，公式应退化回上一问的简单情形

\section{模型的评价与推广}
\subsection{模型的优点}
\subsection{模型的缺点}
\subsection{模型的推广}

% ============================================================================
%  AI 工具使用声明——必须放在【参考文献之前】，二者择一，照抄官方原句
%  依据《全国大学生数学建模竞赛人工智能工具使用规定（2026年试行）》第 3 条
% ============================================================================
\section*{AI 工具使用声明}

% 【情形一】未使用 AI 工具，则保留下面这一句，删除情形二：
本参赛队在竞赛过程中未使用任何 AI 工具。

% 【情形二】使用了 AI 工具，则改用下面这一句（删除情形一），
%           并在支撑材料中放入文件名固定为“AI工具使用详情.pdf”的说明：
% 本参赛队在竞赛过程中使用了 AI 工具，主要用于【简要用途，如语言润色、代码调试等】，
% 详细使用情况见支撑材料。

\begin{thebibliography}{99}
\bibitem{ref1} 作者. 书名或文章名. 出版地: 出版社, 年份.
\bibitem{ref2} 国家标准编号. 标准名称. 年份.
% 工程/工业类题目先查有没有对口国标——两年两次出现：
%   GB/T 2828.1-2012 / GB/T 2828.4-2008（计数抽样检验）
%   GB 12327-1998（海道测量规范）
% 2026 起：参考文献须含网上资料与 AI 工具
\end{thebibliography}

\newpage
% ============================================================================
%  附录：页数不限，但【必须包含完整可运行的源程序】
%  规范第五条：缺程序、程序跑不了、结果与论文不符，都可能取消评奖资格
%  确实没用程序，必须写明“本论文没有用到程序”
% ============================================================================
\appendix
% ctexart 的 \\appendix 默认把标题编成 A、B，全文不出现"附录"二字，
% 评委看不到附录标识。下面这行把编号改成"附录A""附录B"。
\ctexset{section/name={附录,}}
\section{支撑材料文件列表}
\begin{enumerate}[label=(\arabic*),itemsep=2pt]
  \item \texttt{q1\_solve.py} —— 问题一求解主程序
  \item \texttt{result1.xlsx} —— 问题一结果文件
  \item \texttt{AI工具使用详情.pdf} —— 如使用了 AI 工具
\end{enumerate}

\section{源程序}
\subsection{问题一：\texttt{q1\_solve.py}}
\begin{lstlisting}[language=Python]
import numpy as np

def solve():
    """问题一求解主程序。"""
    pass
\end{lstlisting}

\subsection{问题二：\texttt{q2\_optimize.m}}
\begin{lstlisting}[language=Matlab]
% 问题二：最优参数搜索
clear; clc;
\end{lstlisting}

\end{document}
